\chapter{Conclusiones y Trabajo Futuro}
\label{cap:conclusiones}

El objetivo de este estudio consistía en la creación de una herramienta que, dadas unas tareas y unos condicionales, fuese capaz de generar árboles de comportamiento para personajes no jugador que sirviesen como base para rebajar la carga de trabajo de los diseñadores de videojuegos.

Dados los resultados anteriormente mencionados se han llegado a las siguientes conclusiones.

\section{Conclusiones}

% Aunque de forma preliminar debido a la escasez de usuarios encontrados para las pruebas de usuario, se ha llegado a la conclusión de que este modelo es capaz de generar árboles de comportamiento funcionales en un tiempo inferior al de un humano, consiguiendo así el objetivo general planteado.
% Además se ha conseguido implementar un sistema de validación automático que permite reforzar el sistema y que el diseñador pueda observar en tiempo real la ejecución del comportamiento diseñado.
Este trabajo demuestra la viabilidad de generar árboles de comportamiento (BTs) funcionales para NPCs en Unreal Engine 5 a partir de instrucciones en lenguaje natural.
Mediante una arquitectura híbrida que combina Mistral Large 3 para la estructura lógica y llama3:8B para la asignación de variables de la Blackboard, el sistema logra mitigar las alucinaciones y traducir con éxito el pseudocódigo generado a formatos JSON legibles.

Para asegurar la calidad del comportamiento, se implementó un sistema de validación automática doble (un autómata NFA para comprobar dependencias y un validador de objetivos finales).
Las pruebas experimentales con diseñadores (N = 3) confirmaron la alta usabilidad de la herramienta, alcanzando una excelente puntuación SUS de 90.83/100 y demostrando una reducción drástica de los tiempos de desarrollo frente a la creación manual.

Aunque los árboles automáticos presentan una calidad cualitativa inferior debido a la inserción de algunos nodos redundantes, su ejecución en el juego cumple con lo esperado por los usuarios.
La herramienta se consolida como un excelente punto de partida que facilita la edición posterior y disminuye significativamente la carga de trabajo del diseñador.
A pesar de que los resultados cuantitativos son exploratorios debido al tamaño de la muestra, el proyecto valida con éxito la integración de LLMs en flujos de producción de videojuegos.


\section{Trabajo Futuro}
Dadas las limitaciones planteadas en la sección \ref{sec:Limitaciones}, el trabajo a futuro de este estudio se puede centrar en varios esfuerzos:
\begin{itemize}
    \item Con capacidades computacionales mayores estudiar el uso de modelos locales de mayor tamaño para los modelos usados, así como el uso de técnicas de Fine-tuning para comprobar si su uso alcanzaría o mejoraría los resultados actuales con el fin de no depender de APIs de terceros.
    \item Mejorar el sistema RAG con árboles de comportamiento de videojuegos y, adicionalmente, que se guarden en disco el almacén vectorial y se puedan buscar por metadatos gracias al uso de Qdrant, suponiendo así una posible mejora en el contexto dado a los modelos y reduciendo la latencia.
    \item Implementar un agente LLM con una interfaz visual para cambiar los árboles generados de forma más sencilla y que se le pueda dar feedback al generador mediante uso de lenguaje natural.
    \item Implementar un agente LLM que pueda leer las acciones y descripciones de los proyectos para automatizar su suscripción a la Abstract Factory y su descripción para el paso a la herramienta.
    \item Implementar un agente LLM que pueda automatizar las pruebas de validación generando el órden de dependencias y una funciñon de objetivos.
    \item Recopilar más usuarios para las pruebas con el fin de conseguir un número estadísticamente significativo para concluir definitivamente las hipótesis propuestas.
\end{itemize}